\section{Finite certificates and source reconstruction}\label{app:certificates}
This appendix records the algebra used in the main proof, rather than the exploratory calculations that preceded it. In the state order of Section~\ref{sec:model}, the chemical action on a column vector is
\begin{align*}
 (Q_ex)_0&=(x_1-x_0)/50+(x_3-x_0)/50,\\
 (Q_ex)_1&=(51/100)(x_2-x_1)+(1/100)(x_4-x_1)+e(x_0-x_1),\\
 (Q_ex)_2&=2e(x_1-x_2),\\
 (Q_ex)_3&=(51/100)(x_5-x_3)+(1/100)(x_4-x_3)+e(x_0-x_3),\\
 (Q_ex)_4&=(e+1/4)(x_1+x_3-2x_4),\\
 (Q_ex)_5&=2e(x_3-x_5).
\end{align*}
The two-daughter expectation from~\eqref{eq:pairs} is
\[
 Lx=\left(2x_0,\ x_0+x_1,\ \frac{x_0+x_2}{2}+x_1,\ x_0+x_3,\
 \frac{x_0+x_4+x_1+x_3}{2},\ \frac{x_0+x_5}{2}+x_3\right)^{\!\top},
\]
so that $Lw=(28,25,23,98,76,289/2)$ for the weight~\eqref{eq:w}. Combining these with~\eqref{eq:A} gives the exact row ratios below. The last column is the bracket in~\eqref{eq:error} divided by $w_i$, evaluated at $e=.31$; it is maximal there because every rate entering it is nondecreasing in $e$.
\begin{center}\small
\begin{tabular}{@{}lrrrr@{}}
\toprule
Type & $(A_{.01}w)_i/w_i$ & $(A_{.29}w)_i/w_i$ & $(A_{.31}w)_i/w_i$ & Error envelope / $w_i$\\
\midrule
UU & $-\frac{73}{700}$ & $-\frac{73}{700}$ & $-\frac{73}{700}$ & $\frac{353}{700}$\\[5pt]
UR & $-\frac{103}{550}$ & $-\frac{61}{550}$ & $-\frac{29}{275}$ & $\frac{323}{550}$\\[5pt]
RR & $-\frac{21}{125}$ & $-\frac{14}{125}$ & $-\frac{27}{250}$ & $\frac{123}{250}$\\[5pt]
AU & $\frac{559}{4200}$ & $-\frac{421}{4200}$ & $-\frac{491}{4200}$ & $\frac{451}{1050}$\\[5pt]
AR & $-\frac{363}{2150}$ & $-\frac{237}{2150}$ & $-\frac{114}{1075}$ & $\frac{1427}{1075}$\\[5pt]
AA & $\frac{111}{5350}$ & $-\frac{533}{5350}$ & $-\frac{579}{5350}$ & $\frac{477}{2675}$\\[5pt]
\bottomrule
\end{tabular}\end{center}


Each displayed ratio is affine in $e$, so the endpoint values above certify the inequalities~\eqref{eq:nominal} and~\eqref{eq:error} on the whole intervals. Reading off column maxima,
\[
 \max_i\frac{(A_{.01}w)_i}{w_i}=\frac{559}{4200}<\frac{67}{500},\qquad
 \max_i\frac{(A_{.29}w)_i}{w_i}=-\frac{533}{5350}<-\frac{99}{1000},
\]
and the envelope maximum $1427/1075<27/20$ gives the perturbation term $27/2000$.

\paragraph{Reserve inequality.}
The certificate of Theorem~\ref{thm:main} is the single finite rational product
\[
 \frac{278\cdot61\cdot120}{200}\prod_{j=0}^{77}\frac{122}{200-j}<\frac7{10^6},
\]
whose decimal display is $6.9381403704212035\times10^{-6}$; the complete reduced rational value is supplied in the data package. Retaining the denominator $D_{278}=2.5980909040121194\ldots$ gives the sharper value $2.6704763715\ldots\times10^{-6}$ used in Corollary~\ref{cor:filling}, whose initial scale terms are
\[
 \frac{1}{D_{278}}\sum_{j\ge h_0}P_j=
 \begin{cases}
 9.8243173\times10^{-3},&h_0=209,\\
 7.9237537\times10^{-5},&h_0=221,\\
 2.3748386\times10^{-7},&h_0=240,
 \end{cases}
\]
and whose $h_0=208$ value exceeds $10^{-2}$. For Corollary~\ref{cor:smaller}, $D_{161}=2.6271217287\ldots$ at $K=232$ and the sharpened bound is $9.7801264\ldots\times10^{-3}$, against $2.5693583\ldots\times10^{-2}$ for the cruder product.

\paragraph{Exponential inequalities.}
The two exponential bounds used for the delivered wide-box witness follow from the positive partial sums
\[
 \sum_{j=0}^{15}\frac{(99/25)^j}{j!}>50,\qquad
 \sum_{j=0}^{30}\frac{(43/5)^j}{j!}>5000 .
\]
All terms are positive and each sum is less than the corresponding exponential. The exponent identity and marginal assembly are
\[
 108\left(\frac{99}{1000}-\frac{27}{2000}\right)
 -4\left(\frac{67}{500}+\frac{27}{2000}\right)=\frac{2161}{250},
 \qquad
 1-\frac{107}{12500}-\frac7{10^6}=\frac{991433}{10^6}.
\]

\paragraph{The benchmark rows.}
The narrow-box comparison in Table~\ref{tab:examples} uses $A_ew\le-.09w$ on $[.29,.31]$, $A_ew\le.14w$ on $[.01,.31]$, and weighted perturbation at most $.001w$. Its net exponent is $108(.089)-4(.141)=9.048>9$, while $e^9>8000$ by the positive Taylor sum through degree~$25$; its target upper bound is therefore $(214/5)/8000=.00535$. Its healthy bounds come directly from the top-anchor product at $K=M=20$,
\[
 20\cdot\frac{61}{100}\cdot120\cdot\frac{(20(61/100)/\reff)^{10}}{10!}
 \quad\text{at }\reff=12\text{ and }\reff=9,
\]
which give $4.7595\times10^{-4}$ and $8.4518\times10^{-3}$ respectively. By Proposition~\ref{prop:highrenewal} this top-anchor expression is asymptotically exact in $\reff$, so the factor $(12/9)^{10}\approx17.8$ separating the two rows is not an artifact of the bound: at $d=10$ the benchmark genuinely depends on renewal through its tenth power, which is why so much renewal is needed when the reserve carries no spare capacity.
